% Pista alg-23j-q4 · Àlgebra
% Procedència: PAU Catalunya 2023 · Convocatòria de juny · Sèrie 1 · Qüestió 4
\documentclass[11pt,a4paper]{article}
\usepackage{pista}
\pistainfo{Catalunya 2023}{Convocatòria de juny}{Sèrie 1}

\begin{document}

\section*{\textcolor{blauPista}{Qüestió 4}}

\noindent Sigui el sistema d'equacions lineals següent, que depèn del paràmetre real $\lambda$:
\[
\left.\begin{array}{r}
x + 2\lambda y + (2 + \lambda) z = 0 \\
(2 + \lambda) x + y + 2\lambda z = 3 \\
2\lambda x + (2 + \lambda) y + z = -3
\end{array}\right\}
\]

\subsection*{\textit{a)} \normalfont Discutiu el sistema per als diferents valors del paràmetre $\lambda$. \hfill\textnormal{[1,25~punts]}}

\pista{Discuteix el sistema amb el teorema de Rouché-Frobenius:
\begin{itemize}
\item escriu la matriu de coeficients $A$ i la matriu ampliada $A^{'}$;
\item calcula $\det(A)$ i troba els valors de $\lambda$ que l'anul·len;
\item per a cadascun d'aquests valors, compara els rangs d'$A$ i d'$A^{'}$;
\item què pots dir del sistema quan $\det(A) \neq 0$?.
\end{itemize}

\medskip
\textbf{Drecera:} les tres columnes sumen el mateix. Si fas $F_{3} \to F_{1} + F_{2} + F_{3}$ (el determinant no canvia), podràs treure un factor comú de la nova tercera fila.

\medskip
\textbf{Atenció:} comprova si el factor que et queda té arrels reals.}

\subsection*{\textit{b)} \normalfont Per al cas $\lambda = -1$, resoleu el sistema, interpreteu-lo geomètricament i identifiqueu-ne la solució. \hfill\textnormal{[1,25~punts]}}

\pista{Aquest apartat et demana tres coses en una sola frase: «resoleu», «interpreteu geomètricament» i «identifiqueu-ne la solució».
\begin{itemize}
\item \textit{Resoldre.} A l'apartat \textit{a)} ja has vist de quin tipus és el sistema per a $\lambda = -1$: indica-ho. Si és compatible indeterminat, queda't amb dues equacions independents, pren una incògnita com a paràmetre (per exemple $z = t$) i expressa les altres dues en funció de $t$.
\item \textit{Interpretar.} Cada equació és un pla de l'espai: quina posició relativa tenen els tres plans en el cas que t'ha sortit (poden coincidir tots, tallar-se en una recta, no tenir cap punt comú\ldots)?
\item \textit{Identificar la solució.} Separa-la en la part que no depèn de $t$ i la que sí, $(x,y,z) = (\ldots) + t\,(\ldots)$, i interpreta què representa cada part.
\end{itemize}}

\end{document}
